% !TEX root = ../documentation.tex
\chapter{Rule Environments}\label{C:ThemeRuleEnvs}

	Rule environments \RpgPage[p]{S:RuleEnvs} are the elements of a document which depend on the rules of a particular game, such as items, spells and monsters. \rpgtex{} itself only creates them: the theme decides which properties each one accepts, and how it is displayed. \rpgtex{} creates four rule environments, \texttt{RpgItem}, \texttt{RpgSpell}, \texttt{RpgFeat} and \texttt{RpgStat}, each of which may be displayed either as text or as a card. Designers may also create their own.

	Defining a rule environment has three steps:
	\begin{enumerate}
		\item \textbf{Create} the environment, with \cmd{RpgThemeNewRuleEnv} or \cmd{RpgThemeNewCardableRuleEnv}. This is only necessary for new environments, since the four standard ones already exist.
		\item \textbf{Add properties} to it, with the \cmd{Rpg[X]AddProperty} and \cmd{Rpg[X]AddBoolean} commands generated for each environment.
		\item \textbf{Set its format}, with the \cmd{RpgTheme[X]Format} (and, for a card, \cmd{RpgTheme[X]CardFormat}) command generated for each environment.
	\end{enumerate}
	Throughout this chapter, \texttt{[X]} stands for the name of the environment, so that the commands for \texttt{RpgItem} are \cmd*{RpgItemAddProperty}, \cmd*{RpgThemeItemFormat} and so on.

	%% Create the potion used by the examples in this chapter outside of an example, so that it persists
	\RpgThemeNewCardableRuleEnv{Potion}{UsePotionCard}
	\RpgGetExample{des-rule-env}

	\section{Creating Rule Environments}\label{S:ThemeNewRuleEnv}

		\RpgFunction{RpgThemeNewCardableRuleEnv}{[<prefix>]\param{name}\param{switch}[<card-options>][<properties>]}
		{
			Creates the environment \texttt{<prefix><name>} (the prefix is \texttt{Rpg} by default), which may be displayed either as text or as a card, together with a new switch \RpgPage[p]{S:Switches} which chooses between them. The switch is off (text mode) by default. The environment is used as:
			\begin{Code}
				\cmd*{begin}\{<prefix><name>\}[<card-options>]\{<title>\}[<properties>]

				~~<body>

				\cmd*{end}\{<prefix><name>\}
			\end{Code}
			In card mode, the \texttt{card-options} are passed to \envref{RpgCard}; in text mode they are ignored. The optional \texttt{card-options} and \texttt{properties} given to \cmd{RpgThemeNewCardableRuleEnv} are the defaults used when the corresponding argument is omitted. They are replaced, rather than combined, when the argument is given.

			In addition to the commands generated by \cmd{RpgThemeNewRuleEnv}, this creates \cmd{RpgTheme[X]CardFormat}.

			Creating an environment which already exists redefines it, resetting its formats to their defaults; an existing switch keeps its current value.
		}
		\RpgFunction{RpgThemeNewRuleEnv}{[<prefix>]\param{name}}
		{
			Creates the environment \texttt{<prefix><name>} (the prefix is \texttt{Rpg} by default), which is always displayed as text. The environment is used as:
			\begin{Code}
				\cmd*{begin}\{<prefix><name>\}\{<title>\}[<properties>]

				~~<body>

				\cmd*{end}\{<prefix><name>\}
			\end{Code}
			It also creates the commands \cmd{Rpg[X]AddProperty}, \cmd{Rpg[X]AddBoolean} and \cmd{RpgTheme[X]Format}, using the same prefix.

			Creating an environment which already exists redefines it, resetting its format to the default.
		}

	\section{Properties}\label{S:RuleProperties}

		Each property is a key which may be given in the \texttt{properties} argument of the environment. Its value is stored in a variable, which the format can then use. Properties are set separately for each use of the environment, so a property which is not given takes its default value.

		Properties belong to the environment, not to the theme which added them. Adding a property with the same name again replaces it, but a property added by one theme remains available after another theme is loaded.

		\RpgFunction{Rpg[X]AddBoolean}{(*)\param{key}[<aliases>]\param{variable}}
		{
			Adds a flag property, which takes no value. The unstarred form sets the boolean \texttt{variable} to false, and makes it true when the flag is given; the starred form does the opposite. The optional \texttt{aliases} are a comma-separated list of alternative names for the same flag.
		}
		\RpgFunction{Rpg[X]AddProperty}{\param{key}[<aliases>]\param{variable}\param{default}}
		{
			Adds a property which takes a value, stored in \texttt{variable}, with the given \texttt{default}. The optional \texttt{aliases} are a comma-separated list of alternative names for the same property.
		}
		Within either format, the value of a property may also be retrieved by its name with \cmdref{RpgGetKey}. This does not work for flag properties, whose variables should be tested directly (for example, with \cmd*{bool_if:NTF}).

	\section{Formats}\label{S:RuleFormats}

		The format of a rule environment is a code setter (see the sidebar \RpgPage[p]{S:ThemeCommands}). Both formats have the same arguments:
		\begin{RpgTable}{lX}
			Argument & Value
			\\
			\texttt{\#1} & The title of the element.
			\\
			\texttt{\#2} & The body of the environment.
		\end{RpgTable}
		The values of the properties are available through their variables (or \cmd{RpgGetKey}). The following examples restyle the \texttt{RpgPotion} environment created at the start of this chapter; a format set within an example applies only to that example.

		\vbox{
			\RpgFunction{RpgTheme[X]CardFormat}{\param{code}}
			{
				Stores the code which typesets the element in card mode. It is typeset inside an \envref{RpgCard}, and so continues onto further cards if it is too long for one. By default, the title is typeset as a \cmd{subsection}, followed by the body.

			}
			\RpgGetExample{des-potion-card-format}
		}
		\vbox{
			\RpgFunction{RpgTheme[X]Format}{\param{code}}
			{
				Stores the code which typesets the element in text mode. By default, the title is typeset as a \cmd{subsection}, followed by the body.

			}
			\RpgGetExample{des-potion-format}
		}

